\section*{Part VI: The Education Moat: What an Institute Can Build That a Course Seller Cannot}
\addcontentsline{toc}{section}{Part VI: The Education Moat: What an Institute Can Build That a Course Seller Cannot}
\label{sec:part6_the_education_moat}

\subsection{Education Has a Peculiar Economic Problem}
\label{sec:education_peculiar_economic_problem}

Before purchase, learners cannot reliably observe the most important outcome: what they will become capable of doing. Even afterward, they may confuse enjoyment, confidence, a certificate, or career progress with the course's causal impact. Many educational purchases are high stakes, infrequent, socially influenced, and difficult to reverse. Employers and families may value different outcomes than students do. Quality depends partly on learner effort, prior preparation, peer effects, and labor-market opportunity.

This makes education both \textbf{a trust-intensive industry} and \textbf{a measurement-intensive opportunity}. A provider that repeatedly proves real capability gains, communicates them honestly, and delivers them predictably may eventually accumulate an advantage harder to imitate than course content. But a school can also manufacture superficial prestige and misleading placement statistics. The distinction is central to ethical strategy.

\subsection{The Education Product Has 4 Layers}
\label{sec:education_product_four_layers}

\begin{itemize}
    \item \textbf{Layer 1: Information.} What the learner needs to know. Articles, videos, slides, and examples are increasingly abundant and cheap.
    \item \textbf{Layer 2: Practice architecture.} What the learner repeatedly has to do: progressively challenging tasks, opportunities to retrieve knowledge, projects, deadlines, and external feedback.
    \item \textbf{Layer 3: Capability verification.} What independent evidence establishes that they can do it, later and in an unfamiliar setting, rather than merely follow an example.
    \item \textbf{Layer 4: Opportunity and identity.} How learning is translated into credible signals, future access, professional relationships, confidence grounded in competence, and continuing practice.
\end{itemize}

A new institute can begin with a superb Layer 2 and honest Layer 3 even without the prestige needed for a strong Layer 4. Over time, these layers may reinforce 1 another. An institute whose only durable asset is Layer 1 is exposed to free content and AI substitutes.

\subsection{The Core Educational Advantage Is Not the Syllabus}
\label{sec:core_educational_advantage_not_syllabus}

A syllabus describes what will be covered. It tells little about the mental state of the learner after instruction. 2 programs can share module titles (Python, statistics, transformers, deployment) and produce dramatically different capability because of admission criteria, diagnostic work, time on task, assessment, feedback, project authenticity, and motivation.

The hard-to-copy object is \textbf{an integrated transformation system}: a reliable sequence of experiences through which a defined learner changes from 1 demonstrable capability state to another. An effective transformation system makes a testable causal claim, while recognizing that education outcomes are partly co-produced by students.

\subsection{The 1st Education Moat: A Trusted Capability Standard}
\label{sec:first_education_moat_capability_standard}

An institution can become known for what it refuses to certify. Design competency standards around meaningful tasks: debugging a failing model, designing an evaluation, understanding statistical uncertainty, building a reliable service, arguing about costs and security, explaining trade-offs to a colleague. Make criteria observable, calibrated, and consistently applied.

If a certificate means ``paid fees and watched lessons,'' it is a weak signal. If it means ``passed unseen, individually authenticated tasks against a published standard,'' it could become more credible, \textbf{but only after external parties independently come to trust the process and results}. A self-created badge is not automatically labor-market currency.

\subsection{Trustworthy Assessment Beats Prestigious Vocabulary}
\label{sec:trustworthy_assessment_beats_prestige_language}

A school that claims ``industry-grade assessment'' should establish what the tasks predict. Are they representative of actual work? Are scores reliable across graders? Do they distinguish prepared from unprepared candidates? Can an excellent memorizer pass without being competent? Does performance persist after a delay? How much does an AI assistant change the meaning of the score?

The assessment bank itself may be copied or leaked; the deeper capability is \textbf{a living process of task design, calibration, security, evidence review, and updating}. Use multiple task versions, unseen transfer tasks, clear permitted-tool policies, and recorded oral defenses when warranted, subject to privacy safeguards.

\subsection{A Learning-Outcome Evidence Hierarchy}
\label{sec:learning_outcome_evidence_hierarchy}

\begin{table*}[t]
\centering
\small
\caption{The Learning-Outcome Evidence Hierarchy: Triangulating Educational Impact}
\label{tab:learning_outcome_hierarchy}
\begin{tabularx}{\textwidth}{p{4cm} X X}
\toprule
\textbf{Evidence Level} & \textbf{What It Establishes} & \textbf{What It Does Not Establish} \\
\midrule
Attendance and completion & Participation, compliance, and persistence & Independent mastery or transfer \\
\addlinespace[2pt]
Positive feedback surveys & Student satisfaction and perceived usefulness & Objective learning gains \\
\addlinespace[2pt]
Portfolio project artifact & Ability to produce an artifact under some guidance & Independent authorship or unseen transfer \\
\addlinespace[2pt]
Proctored / new task & Performance on a specific independent task & Retention durability and broader domain transfer \\
\addlinespace[2pt]
Baseline-to-endline delta & Observed improvement in a measured domain & Program causality without matched controls \\
\addlinespace[2pt]
Delayed unseen assessment & Longitudinal retention and some conceptual transfer & Performance under messy workplace conditions \\
\addlinespace[2pt]
Employer-verified impact & Practical workplace relevance in a specific context & Direct causal attribution solely to school \\
\addlinespace[2pt]
Credible comparison design & Strong causal evidence of institutional contribution & Guarantee of identical success for every future student \\
\bottomrule
\end{tabularx}
\end{table*}

Good evidence accumulates through triangulation. Do not combine unlike measures into a single ``success rate'' without transparent definitions. Publish denominators, attrition, missing data, and consent-based case detail.

\subsection{The Problem of Selection and Causal Attribution}
\label{sec:selection_and_causal_attribution}

A school that enrolls already exceptional students will have impressive alumni even if it adds relatively little. A school serving beginners may produce greater gains while showing lower final scores. Peer effects and learner motivation can improve outcomes, but such improvements complicate causal claims.

The general lesson of value-added research is not that educational impact is impossible to measure; it is that \textbf{baseline, case mix, and causal design matter}. For a small Bayes cohort, begin with diagnostics, standardized rubrics, delayed tasks, and transparent longitudinal tracking. Avoid claiming causal educational superiority on small uncontrolled samples. As scale and resources grow, consider matched comparisons, randomized access when ethically and operationally appropriate, or independent program evaluation.

\subsection{The 2nd Moat: Accumulated Misconception Intelligence}
\label{sec:second_moat_misconception_intelligence}

Technical educators often notice predictable failure modes. A learner can reproduce a backpropagation derivation but misunderstand when a gradient vanishes. A developer can call an LLM API but fail to evaluate retrieval quality or handle prompt injection. A modeler can achieve strong backtest results while leaking future information.

If instructors record these mistakes in a structured way (task context, prior skill, misconception, intervention, and later result), the institution can accumulate a \textbf{map of where learners fail and which interventions help}. Over time this can guide better diagnostics, feedback, and sequencing. The moat is the validated \textit{response system}, not a pile of error logs.

\subsection{The 3rd Moat: Reproducible Faculty Excellence}
\label{sec:third_moat_reproducible_faculty}

Every learner values excellent teachers; most schools struggle to make that excellence consistent. A founder might have unusually deep subject knowledge and an intuitive style, but the institution only compounds when it can hire, develop, and calibrate other educators.

Build a teaching-faculty system with: realistic job previews; subject-matter work samples; teaching auditions with mixed-ability learners; live diagnosis of misconceptions; feedback-quality assessments; structured observations; review of learner work; mentorship for new instructors; explicit performance standards; and periodic calibration. A perfect audition is not enough: evaluate actual learner progress and student experience across time, with reasonable controls.

\subsection{Faculty Quality and the Instructor-Scarcity Paradox}
\label{sec:faculty_scarcity_paradox}

An institute may be strengthened by hiring exceptional teachers and weakened by depending entirely on them. A star faculty member can create demand, but departure, compensation negotiation, or scheduling conflict exposes the business. The same applies to a founder whose personal credibility supplies most sales.

The answer is not to make people interchangeable. It is to \textbf{make individual excellence more productive and more transferable}: co-teaching, shared task banks, peer observation, critique sessions, lesson studies, common rubrics, documented interventions, and the option for promising graduates to become teaching assistants under supervision. Faculty should have reasons to stay: meaningful work, growth, fair pay, autonomy, intellectual community, not artificial restrictions.

\subsection{Do Not Build the Faculty Moat by Underpaying Teachers}
\label{sec:dont_underpay_teachers}

A model that promises premium results while paying below-market wages may temporarily enjoy attractive margins but will struggle to recruit and retain capable people. Faculty development requires training time and supervision. A sustainable education moat must align tuition, class size, instructor compensation, scholarship policy, assessment intensity, and financial reserves.

\textbf{The relevant efficiency measure is not learner-to-teacher ratio alone; it is verified capability gained per unit of scarce expert attention.} Use software, peers, teaching assistants, and well-designed asynchronous materials to reserve experts for moments where judgment matters most.

\subsection{The 4th Moat: A Distinctive Peer Environment}
\label{sec:fourth_moat_peer_environment}

Peers can accelerate learning through comparison, explanation, collaboration, accountability, norms, and access to diverse perspectives. They can also discourage beginners, spread weak practices, or facilitate plagiarism. Selectivity can help form a serious learning environment, but selection based only on prior achievement can sacrifice access and confuse institutional contribution with intake quality.

For Bayes, favor observable curiosity, commitment, baseline readiness, willingness to work, and integrity rather than status alone. Design peer value intentionally: code reviews, structured critiques, demonstrations, study groups, pair debugging, and alumni technical sessions. Measure \textit{interaction quality}, not Discord membership counts.

\subsection{Are Cohorts Network Effects?}
\label{sec:are_cohorts_network_effects}

Sometimes, but usually not automatically. If a 30-person cohort offers more diverse project feedback than a 15-person cohort, added participants might improve value. If a 300-person cohort overwhelms feedback and reduces trust, the marginal participant could reduce value.

A peer network becomes more defensible when members repeatedly provide \textbf{valuable, difficult-to-replace interaction}, with dense participation and norms that support quality. An alumni group of 5,000 with 5 active people may be economically weaker than a highly engaged network of 150 credible professionals.

\subsection{The 5th Moat: An Alumni Relationship That Persists}
\label{sec:fifth_moat_persisting_alumni_relationship}

Education providers often define customer success as graduation. But technical fields change for decades. An institute can remain valuable through alumni seminars, updated learning paths, reviewed practice challenges, career-long learning records, specialist communities, mentoring, and pathways to teach or collaborate.

This relationship can create re-enrollment, referral, employer access, reputation, and shared institutional identity. However, \textbf{alumni membership is a stock; alumni interaction is the flow that sustains its value}. Design privacy, membership expectations, portability, and consent clearly. Do not use alumni identities as marketing assets without permission.

\subsection{The 6th Moat: Genuine Employer Relationships}
\label{sec:sixth_moat_genuine_employer_relationships}

A credible employer connection is not a list of HR logos on a webpage. It means repeated interaction that makes both sides better off: the institution understands real skill requirements; employers trust the assessment evidence; students receive authentic problems or appropriate opportunities; and employers see enough value to return.

Begin with modest commitments: guest critique, technical review, authentic project briefs, anonymized skill-gap conversations, or paid upskilling pilots, not promises of placement. The institution must be careful about conflicts of interest, student consent, confidential company data, and the temptation to design the curriculum only for the current hiring fad.

\subsection{Employer Network Effects: Attractive but Often Weak Initially}
\label{sec:employer_network_effects_weak_initially}

More skilled learners can attract more employers; more employers can attract learners. This resembles a 2-sided network effect. Yet employers may already have abundant recruiting channels and may not meaningfully change behavior for a small cohort. Employers may multihome, hire infrequently, or prefer familiar credentials. A school should not assume an employer marketplace exists because 3 people agree to guest lecture.

Measure repeated employer use, time-to-qualified-candidate, quality of matches, hiring or project engagement, and independent employer trust. \textbf{2-sided market power only emerges after each side gains demonstrable value from the other.}

\subsection{The 7th Moat: Referenceable Instructional Quality}
\label{sec:seventh_moat_referenceable_quality}

Excellent delivery creates recommendations; credible recommendations reduce perceived risk; more appropriately matched students improve classroom dynamics; better outcomes create more references. This could become 1 of the strongest loops for a small school.

But selection bias and feedback manipulation are serious dangers. Publish real learner experiences with consent; present both the range and the limitations of outcomes; record refund disputes and complaints; fix problems visibly. A school becomes trusted through \textbf{accurate expectations fulfilled repeatedly}, not only celebratory testimonials.

\subsection{The 8th Moat: A Proprietary Case and Project Pipeline}
\label{sec:eighth_moat_project_pipeline}

Authentic technical work is frequently more educational than generic toy assignments. An institute can form relationships with startups, research labs, enterprises, civic bodies, or open-source maintainers to obtain practical, appropriately scoped problem briefs. That generates relevance, evidence of learner capability, faculty learning, and institutional relationships.

The pipeline must have reliable scoping and quality control. Real clients can change requirements, expose confidential data, or demand production work from inexperienced learners. Keep educational duties primary. Use contracts clarifying rights, confidentiality, supervision, and expectations; do not treat unpaid student labor as an economic loophole.

\subsection{The 9th Moat: A Continuously Updated Capability Map}
\label{sec:ninth_moat_updated_capability_map}

A technical school can map underlying skills (not just technologies) to learning evidence. For example: statistical reasoning, debugging, system decomposition, reliability engineering, experimental design, AI evaluation, safe deployment, technical writing, and decision making under uncertainty.

When tools change, the course can update surface tasks while preserving deeper standards. The institution's strength is knowing \textbf{what capability matters, what indicates it, and how to develop it efficiently}. That is more durable than allegiance to 1 language, model family, or fashionable acronym.

\subsection{The 10th Moat: The Educational Operating System}
\label{sec:tenth_moat_educational_operating_system}

Picture the systems that make instruction reproducible: admission diagnostic; learner readiness grouping; timetable; recorded preparation; structured live problem solving; graded practice; feedback escalation; technical office hours; assessment integrity; faculty calibration; learner-support protocols; and longitudinal analytics.

Any individual component can be copied. What is harder is building a mature operating system where staff execute consistently, data flows to decisions, and leadership closes the feedback loop. The system has to remain humane: rigid uniformity that ignores learners will undermine the product.

\subsection{The 11th Moat: Cost-Efficient Quality Without Dilution}
\label{sec:eleventh_moat_cost_efficient_quality}

Premium education often assumes a high marginal human cost. Scale-sensitive advantages come from designing high-quality repeatable learning components and reserving expert attention for genuine bottlenecks. Use short reusable explanations, self-checks, low-stakes automated practice, AI tutoring with verified guardrails, peer review with calibration, TA escalation, and well-timed expert critique.

\textbf{Do not equate automation with cost advantage until outcomes are held constant.} If automated feedback produces confident misconceptions, it can increase remediation cost. Measure time-to-mastery, verified independent performance, student hours, expert minutes, and total support costs by learner segment.

\subsection{The 12th Moat: Institutional Research and Innovation}
\label{sec:twelfth_moat_institutional_research}

A mature technical institute can investigate learning methods, assessment validity, tooling, and emerging scientific practice. Applied research can improve teaching, attract faculty, strengthen employer collaboration, and eventually create intellectual property or spinouts. But research reputation is built through actual contributions and external scrutiny, not a tab on the website labeled ``Research.''

For Bayes, early research might be modest: publishing a careful case study of how working engineers improve on unseen AI-system evaluation tasks, with limitations made explicit. The long-term ambition of a research-intensive campus is coherent only if revenues, talent, infrastructure, and research standards evolve in that direction.

\subsection{Curriculum Content Is Rarely a Strong Secret Moat}
\label{sec:curriculum_content_rarely_secret_moat}

Hiding PDFs or refusing to disclose program structure is often a poor defense. Students need transparency. Competitors can assemble similar materials from the open internet. Proprietary slides rarely substitute for strong teaching.

Protect actual original materials with appropriate rights while showing enough work publicly to earn trust. Invest the bulk of defensibility effort into \textbf{delivery quality, authentic evaluation, faculty reproduction, relationships, and longitudinal evidence}. Publishing thoughtful teaching can itself become distribution and credibility.

\subsection{Prestige: A Powerful Moat, but Not a Valid Starting Assumption}
\label{sec:prestige_late_stage_moat}

Academic prestige can attract high-performing applicants, faculty, employers, donors, research opportunities, and partners. These inputs can raise subsequent output quality and reinforce prestige. It is a form of accumulated social belief, sometimes supported by real scarce resources and self-selection.

An unknown institute cannot responsibly manufacture the history underlying that belief. It can begin building \textit{earned reputation} in 1 demanding niche, publish honest evidence, and gradually expand. Avoid artificial scarcity, invented advisory boards, borrowed logos, fake affiliations, and rankings without a defensible methodology. Prestige worth having is a late-stage result of institutional conduct.

\subsection{Accreditation and Degrees: Not an Early Shortcut}
\label{sec:accreditation_not_early_shortcut}

Formal credentials can become powerful through legal authority, employer recognition, academic standards, and compatibility with institutional pathways. However, regulatory regimes vary by jurisdiction, type of provider, and proposed award. Recognition cannot be assumed from registering a company or teaching sophisticated material.

A lean non-degree institution can still offer substantial value. Its initially honest claim is: \textbf{we teach and assess specified competencies}, not ``our certificate has the labor-market authority of a university degree.'' A future campus or accredited program is a separate capital, governance, faculty, regulatory, and academic problem.

\subsection{The Quality-Selectivity Trade-Off}
\label{sec:quality_selectivity_tradeoff}

Selective admissions can improve peer environment and observed graduate outcomes. But extreme selectivity can reduce the school's actual educational contribution and exclude people who would gain substantially from teaching. A top-score school that accepts only top-score students may look powerful while doing less transformation work than a school with a more challenging intake.

The right policy depends on the program's promised transformation and practical support capacity. Establish \textbf{readiness thresholds for safe participation}, not a prestige contest. If the cohorts include materially different readiness groups, offer bridges or multiple tracks so the highest standard is not achieved by silently abandoning beginners.

\subsection{The Quality-Access Trade-Off}
\label{sec:quality_access_tradeoff}

A school that charges premium tuition may fund excellent teaching but restrict access. A school that underprices unsustainably may widen entry briefly and then fail. Scholarships can be valuable if financed transparently and not by deceptive cross-subsidies. Enterprise training or college partnerships might help fund more accessible direct cohorts, but only if those verticals actually earn contribution and do not drain instructional capacity.

Access is a strategic design choice, not evidence that premium pricing is immoral or that low pricing is virtuous. The institutional test is \textbf{net social value, fairness, financial sustainability, and the ability to keep promises}.

\subsection{B2C Education and B2B College Partnerships Have Different Moats}
\label{sec:b2c_and_b2b_different_moats}

A direct learner usually evaluates personal transformation, price, time, instructor quality, peer fit, and risk. A college buyer evaluates alignment with curriculum, faculty workload, scheduling, administrative feasibility, student outcomes, procurement risk, and reputational safety. An enterprise L\&D buyer cares about changed workflows, measurable business performance, compliance, and managerial adoption.

These are not simply 3 sales channels for 1 course. They require different delivery systems, contracts, metrics, and relationships. Their underlying capability assets may overlap, but their respective switching costs and repeat dynamics differ. A good strategy allows shared capabilities without assuming 1 sales playbook fits all.

\subsection{When Verticals Reinforce 1 Another, and When They Do Not}
\label{sec:when_verticals_reinforce_one_another}

A direct cohort may generate instructional innovations valuable to partner colleges. Employer projects may provide authentic cases useful across cohorts. College faculty training may improve Bayes's instructor-development routines. Enterprise AI deployments may reveal urgent industry skills.

That is a plausible \textbf{knowledge flywheel}. It is not guaranteed: enterprise consulting can consume founder hours, pull curriculum toward specific proprietary stacks, and cause delivery delays. College partnerships may prioritize seat volume over instructional rigor. Build shared asset governance and a clear hierarchy of promises so growth in 1 vertical cannot degrade another.

\subsection{A Common Moat Illusion: The Placement Network}
\label{sec:common_illusion_placement_network}

A provider may claim a placement moat based on recruiters, yet employers have many alternatives and the school may have little influence on hiring decisions. The hardest part may be labor-market demand, not contact introductions. Unverified placement numbers can conceal selection, missing outcomes, candidate experience, or salary definitions.

Build authentic employer trust and verifiable technical evidence; never guarantee employment outcomes you do not control. Distinguish \textbf{employer introductions}, \textbf{interviews}, \textbf{offers}, \textbf{accepted jobs}, \textbf{salary changes}, and \textbf{causal program contribution}. These are different measures.

\subsection{Another Illusion: A Huge Library of Content}
\label{sec:illusion_huge_content_library}

If thousands of videos are inexpensive to generate, the marginal value of an additional recorded lesson falls unless it addresses a precise learning need better than alternatives. Large libraries may create choice overload, weak completion, and stale material.

Content becomes more defensible when integrated into a diagnosis-driven learning pathway with carefully designed practice and verified feedback. \textbf{Curriculum coherence and instructional judgment} often matter more than raw course count.

\subsection{Another Illusion: A Giant Alumni Group}
\label{sec:illusion_giant_alumni_group}

An alumni chat group is a communication channel, not necessarily a network advantage. The relevant metrics are repeated useful interactions, diversity of useful expertise, peer trust, quality of mentoring, genuine professional opportunities, and propensity to renew relationships.

If alumni primarily receive promotional messages, membership can become a negative signal. Design specific rituals that create reciprocal value: quarterly technical seminars, reproducible project critiques, mentorship matching with clear expectations, and carefully facilitated problem-solving circles.

\subsection{Another Illusion: ``AI-Powered Personalization''}
\label{sec:illusion_ai_powered_personalization}

Many providers can buy access to similar AI tools. Personalization is valuable only when diagnosis is accurate, interventions are effective, learner data are used with appropriate permission, human oversight resolves consequential errors, and the improvement is demonstrated. Generating a custom learning plan is not the same as changing a person's capability.

A meaningful advantage may lie in \textbf{which misconception is targeted, when, with what feedback, and how success is checked}. The user interface around those decisions may remain simple.

\subsection{Another Illusion: Founder Celebrity}
\label{sec:illusion_founder_celebrity}

A visible founder may lower acquisition friction and create initial trust. But celebrity can become a single point of failure. If the founder cannot teach, supervise, sell, and care for everyone, growing demand can diminish quality. Moreover, a persona that attracts attention may conflict with the intellectual humility needed for rigorous education.

The better goal is to gradually shift the buyer's trust from \textbf{``I know this founder''} to \textbf{``I know what this institution reliably does.''} Preserve the founder's identity as an amplifier, not as the sole source of value.

\subsection{The Compounding Sequence for an Unknown School}
\label{sec:compounding_sequence_unknown_school}

The sequence is usually not \textbf{brand} $\to$ \textbf{enrollment} $\to$ \textbf{outcomes}. For a credible new provider, it can be:
\begin{center}
\small
\textbf{Founder competence} $\to$ \textbf{small paid delivery} $\to$ \\
\textbf{independently assessed learner progress} $\to$ \textbf{honest evidence} $\to$ \\
\textbf{better-matched applications} $\to$ \textbf{improved teaching routines} $\to$ \\
\textbf{faculty replication} $\to$ \textbf{trusted referrals and repeat partners} $\to$ \\
\textbf{progressively stronger institutional reputation}
\end{center}

Notice where the hardest work happens: between instruction and evidence, and between founder performance and institutional reproduction. Those steps make the loop resistant to superficial copying.

\subsection{Design the Educational Moat Around 1 Narrow Capability Promise}
\label{sec:design_moat_narrow_capability_promise}

At day 0, choose a learner and a transformation you can credibly deliver and verify. ``We teach the future of AI'' is too broad. A more testable hypothesis might be: ``Working software developers can become capable of designing, evaluating, and deploying a narrowly specified category of reliable AI-enabled system, demonstrating the result on unfamiliar tasks.''

The program should not be built to look extensive. It should be built so that \textbf{a buyer can understand what makes success observable}, and the institution can repeatedly make success more likely. As evidence grows, the scope of the promise can expand.

\subsection{The Ethical Limit to an Educational Moat}
\label{sec:ethical_limit_educational_moat}

Education should increase people's capability and agency. A defensibility strategy that depends on hiding information, misleading families, trapping students in financing, producing unearned credentials, exploiting unpaid work, or creating dependence on the institute is strategically dangerous and institutionally contradictory.

The most desirable moat is 1 where \textbf{students become more independent, even while alumni voluntarily remain connected}. A school that becomes less necessary to its graduates while becoming more trusted by them has a stronger claim to long-run legitimacy than 1 that manufactures permanent dependence.
